% =============================================================================
% APPENDICE: Nota metodologica sulla scrittura aumentata (200-350 parole)
% BOZZA da rivedere e confermare dall'autore: i fatti riportati sono ricavati
% da METODOLOGIE.txt e dai rapporti di verifica in revisione/.
% =============================================================================

\chapter*{Nota metodologica: come è stato scritto questo libro}
\addcontentsline{toc}{chapter}{Nota metodologica}

Questo libro è nato da appunti, letture e conversazioni accumulati in alcuni anni. Le tesi, la struttura in quattro parti, la scelta dei casi e delle esperienze personali sono mie. Per la stesura ho usato strumenti di intelligenza artificiale generativa, con un metodo preciso: per ogni capitolo scrivevo prima una traccia di poche centinaia di parole, con i punti da toccare, le fonti principali e gli esempi; poi lavoravo una sezione alla volta, chiedendo al modello proposte di formulazione che rileggevo, correggevo o scartavo. L'IA mi ha aiutato soprattutto nel ritmo della prosa, nella ricerca di collegamenti tra capitoli e, nella fase finale, nella revisione: individuare ripetizioni, controllare i rimandi interni, segnalare affermazioni troppo sicure.

Ogni fonte è stata poi verificata separatamente, citazione per citazione. È stato il passaggio più istruttivo. Le verifiche hanno fatto emergere errori di tipo diverso: articoli scientifici con autori, titolo e rivista plausibili ma inesistenti; un libro attribuito a uno storico che in realtà corrispondeva a una sua conferenza; epigrafi attribuite a studiosi famosi che non le avevano mai scritte; dati corretti nella sostanza ma riportati con troppa sicurezza. Sono esattamente gli errori descritti nel cap.~\ref{cap:p3c11}: testi fluidi e credibili, che la mia mente era pronta ad accettare perché confermavano ciò che volevo dire. Le fonti non verificabili sono state eliminate o sostituite; le frasi senza autore accertato sono diventate mie, e firmate come tali.

Restano certamente errori che non ho visto. Per questo la bibliografia indica, per le affermazioni principali, il grado di evidenza su cui poggiano. Il resto è un invito: se ne trovate, scrivetemi.

\medskip
\noindent\dascrivere{verificare e completare questa nota con i dati reali: periodo di scrittura, strumenti usati (nomi e versioni), modalità concrete di verifica delle fonti, un recapito per le segnalazioni dei lettori.}
